\documentclass[spanish,letterpaper,oneside]{article}
\usepackage[utf8]{inputenc}
\usepackage[T1]{fontenc}
\usepackage[spanish,es-nodecimaldot]{babel}
\usepackage{helvet}
\renewcommand{\familydefault}{\sfdefault}
\usepackage[margin=1.4cm,top=1.55cm,bottom=1.55cm]{geometry}
\usepackage{amsmath,amssymb,booktabs,array,longtable,tabularx,colortbl,tikz,fancyhdr,enumitem,needspace}
\usepackage[most]{tcolorbox}
\usetikzlibrary{arrows.meta,positioning,calc,babel}
\usepackage{xcolor}
\usepackage{hyperref}

\definecolor{iiiepeNavy}{HTML}{163B5C}
\definecolor{iiiepeBlue}{HTML}{1F6F9F}
\definecolor{iiiepeAqua}{HTML}{20B7A8}
\definecolor{iiiepeOrange}{HTML}{F28C28}
\definecolor{iiiepePaper}{HTML}{F6F8F7}
\definecolor{iiiepeInk}{HTML}{1E2A32}
\definecolor{iiiepePale}{HTML}{EEF5F8}
\definecolor{okGreen}{HTML}{1B6B3A}
\definecolor{warnAmber}{HTML}{B25A00}

\hypersetup{colorlinks=true,linkcolor=iiiepeBlue,urlcolor=iiiepeBlue}
\setlength{\parindent}{0pt}
\setlength{\parskip}{5pt}
\renewcommand{\arraystretch}{1.22}
\arrayrulecolor{iiiepeNavy!55}

\pagestyle{fancy}
\fancyhf{}
\fancyhead[L]{\scriptsize\color{iiiepeNavy}\textbf{IIIEPE · Maestría en Aprendizaje y Enseñanza de las Matemáticas}}
\fancyhead[R]{\scriptsize\color{iiiepeNavy}\textbf{PFM02 Temas Selectos de Matemáticas I · SD26}}
\fancyfoot[L]{\scriptsize Mtro. Martín Jonathan de la Cruz Muñoz (Matrícula 00874)}
\fancyfoot[C]{\scriptsize\thepage}
\fancyfoot[R]{\scriptsize Guía de Entregas: Presencial vs. Plataforma}
\renewcommand{\headrulewidth}{0.5pt}

\newtcolorbox{guiabox}[2][]{%
  enhanced,breakable,
  colback=iiiepePaper,colframe=#2,
  boxrule=1pt,leftrule=4pt,arc=2pt,
  left=8pt,right=8pt,top=6pt,bottom=6pt,
  fonttitle=\bfseries\color{white},coltitle=white,
  colbacktitle=#2,#1}

\begin{document}
\color{iiiepeInk}

\begin{center}
  {\Large\bfseries\color{iiiepeNavy} INFORME GUÍA DE ENTREGAS Y ORGANIZACIÓN}\\[4pt]
  {\normalsize\bfseries Ruta clara para la entrega presencial (mañana) vs. plataforma EVA IIIEPE}\\[4pt]
  {\small\textbf{Unidad de Aprendizaje:} PFM02 Temas Selectos de Matemáticas I (Ago--Nov 2026 Grupo A) · \textbf{Docente:} Mtra. Denis Concepción Escobedo Martínez}\\
  {\small\textbf{Estudiante:} Mtro. Martín Jonathan de la Cruz Muñoz (Matrícula \texttt{00874}) · \textbf{Corte:} 20 de noviembre de 2026}
\end{center}

\vspace{2pt}
\begin{guiabox}[title={Lectura rápida en 30 segundos (para no confundir plataforma vs. entrega presencial)}]{iiiepeNavy}
\begin{itemize}[leftmargin=1.4em,itemsep=2pt]
  \item \textbf{¿Qué SÍ está abierto todavía en plataforma?} Solo el buzón \textbf{Portafolio de evidencias Bloque I} (escala $0\text{--}5$, actualmente \textit{Sin envíos}, en caso de que también se requiera digital; si la docente solo lo recibe físico mañana, ya está compilado completo en 39 páginas dentro de \texttt{00\_Portafolio-Evidencias-Bloque-I/}).
  \item \textbf{¿Qué ya está entregado y calificado en plataforma?} \textbf{Actividades 2, 3 y 4} ($70.00/100.00$ cada una) y \textbf{Trabajo Final} (actualizado hoy a las \textbf{19:00} con el video limpio \texttt{4:46}, enlaces YouTube y Drive, \textit{Guión} y reportes Act 5 y Act 6; en pantalla aún muestra la calificación previa de las \textbf{14:45}: $25.00/36.00$). Todos los reportes SD26 están validados editorialmente al nivel \textbf{MAESTRO (EMS 100/100, Estilo 100/100)} con bibliografía formal BibTeX y verificación matemática independiente.
  \item \textbf{¿Qué NO tuvo buzón individual en plataforma y se entrega presencialmente?} \textbf{Actividad 1} (tablero colaborativo Canva Equipo 4), \textbf{Examen Parcial 1} (\textit{Sospechoso, muy sospechoso}), los \textbf{retos impresos} de Actividades 2 y 3, las \textbf{hojas verdes} de Actividad 5 (Guion de la Habilidad + Solución a los retos) y las \textbf{lecturas a mano} de Actividad 4.
\end{itemize}
\end{guiabox}

\section{Mapa maestro de actividades, retos y estado (Plataforma vs. Presencial)}

{\footnotesize\begin{longtable}{|>{\raggedright\arraybackslash}p{0.19\textwidth}|>{\raggedright\arraybackslash}p{0.23\textwidth}|>{\raggedright\arraybackslash}p{0.28\textwidth}|>{\raggedright\arraybackslash}p{0.25\textwidth}|}\hline
\rowcolor{iiiepeNavy}
\color{white}\textbf{Actividad / Evaluable} &
\color{white}\textbf{Estado en plataforma EVA (20-nov-2026)} &
\color{white}\textbf{Qué entregar mañana presencial (Mtra. Denis)} &
\color{white}\textbf{Carpeta y archivo clave en \texttt{entregas-...SD26/}} \\\hline
\endfirsthead\hline
\rowcolor{iiiepeNavy}
\color{white}\textbf{Actividad / Evaluable} &
\color{white}\textbf{Estado en plataforma EVA} &
\color{white}\textbf{Qué entregar mañana presencial} &
\color{white}\textbf{Carpeta y archivo clave} \\\hline
\endhead

\textbf{Portafolio de Evidencias Bloque I} (Periodo anterior: Fracciones, Decimales, Signos, Proporcionalidad y Ecuaciones Lineales) &
Buzón abierto \textit{Sin envíos}, rango $0\text{--}5$. Botón \textit{Agregar entrega} disponible si desea respaldo digital además del físico. &
\textbf{Impreso o USB (prioritario mañana):} Portafolio completo de \textbf{39 páginas} (Portada con índice paginado + los 5 reportes integrados del Bloque I). &
\texttt{00\_Portafolio-Evidencias-Bloque-I/}\\
\texttt{Portafolio-Evidencias-Bloque-I\_}\\
\texttt{Completo\_Martin-Jonathan-de-la-}\\
\texttt{Cruz-Munoz.pdf} (39 pp., 1.43 MB). \\\hline

\textbf{Actividad 1 (Sesión 8):} \textit{¿Qué sabemos y cómo lo enseñamos?} + Retos 1, 2 y 3 &
\textbf{Sin buzón en plataforma.} Evaluación colaborativa en tablero Canva (\textbf{Equipo 4}: Minecraft/Papercraft y jardín $8\times 6$). &
Reporte con los 3 retos de área y perímetro (esquemas TikZ y parábola $72\text{ m}^2$) + evidencias PNG del Equipo 4 en Canva. &
\texttt{01\_Actividad-1\_.../} (\texttt{Actividad 1\_...docx}, \texttt{reporte-...act-01...pdf}, \texttt{evidencia-canva-*.png}). \\\hline

\textbf{Actividad 2 (Sesión 9):} \textit{El lugar perfecto} (Mediatriz, Bisectriz, Circuncentro e Incentro) &
\textbf{Enviado 19-nov · Calificado $70.00/100.00$} (20-nov 14:47). Incluye \textit{Guión 2}, \textit{reporte}, \textit{presentación} y 2 \texttt{.ggb}. &
\textbf{Hojas de retos presenciales} (Reto 1: 2 salones/mediatriz; Reto 2: 2 caminos/bisectriz; Reto 3: 3 salones/circuncentro) con esquemas y GeoGebra. &
\texttt{02\_Actividad-2\_.../}\\
\textbf{\texttt{Retos-Resueltos-Actividad-02\_...pdf}} (3 pp.) + \texttt{reporte-...act-02...pdf} + \texttt{.ggb}. \\\hline

\textbf{Actividad 3 (Sesión 10):} \textit{Universo 3-60} (Triángulos, Polígonos inscritos y Ángulo inscrito) &
\textbf{Enviado 19-nov · Calificado $70.00/100.00$} (20-nov 14:48). Incluye \textit{Guión 3}, \textit{reporte}, \textit{presentación} y \texttt{.ggb}. &
\textbf{4 hojas presenciales resueltas:} Justificación/Demostración $180^\circ$, 2 circunferencias púrpuras ($n=6,8$), ficha de 3 polígonos y tabla de 9 polígonos hasta $n=24$. &
\texttt{03\_Actividad-3\_.../}\\
\textbf{\texttt{Retos-Resueltos-Actividad-03\_...pdf}} (4 pp.) + \texttt{reporte-...act-03...pdf} + \texttt{.ggb}. \\\hline

\textbf{Examen Parcial 1 (Sesión 10):} \textit{Parcial 1. Sospechoso, muy sospechoso} &
\textbf{Sin buzón en plataforma.} Instrumento de evaluación presencial de Sesión 10. &
\textbf{Hoja de examen resuelta} (en físico o transcrita al formato original con puño y letra si así lo pide) + página 2 con validaciones geométricas. &
\texttt{04\_Examen-Parcial-1\_.../}\\
\textbf{\texttt{Parcial-1\_Sospechoso-muy-sospechoso\_...pdf}} (y \texttt{.docx}). \\\hline

\textbf{Actividad 4 (Sesión 11):} \textit{Construyendo aprendizaje a través de habilidades} (SEP 1994 y Alsina 2023) &
\textbf{Enviado 19-nov · Calificado $70.00/100.00$} (20-nov 14:51). Incluye \textit{Guión 4}, \textit{reporte} y \textit{presentación}. &
Confirmar con la profesora el rubro de \textbf{reportes de lectura a mano} (20\,\% del encuadre inicial) llevando el resumen impreso y/o manuscrito de SEP (1994) y Alsina (2023). &
\texttt{05\_Actividad-4\_.../}\\
\texttt{Guion 4\_...docx}, \texttt{reporte-...act-04...pdf}. \\\hline

\textbf{Actividad 5 (Sesión 12):} \textit{¿Cómo hacer visible el desarrollo de habilidades?} (Guion y retos) &
\textbf{Sin buzón propio;} \texttt{reporte-...act-05...pdf} ya quedó adjunto dentro del \textbf{Trabajo Final}. &
\textbf{2 hojas verdes presenciales:} Pág. 1 = \textbf{Guion de la Habilidad} (\textit{Generalizar}); Pág. 2 = \textbf{Solución a los 7 retos + Comunicar} con esquemas TikZ. &
\texttt{06\_Actividad-5\_.../}\\
\textbf{\texttt{Actividad-05\_Guion-Habilidad-y-}}\\
\textbf{\texttt{Retos-Resueltos\_...pdf}} (2 pp.) + \texttt{reporte-...act-05...pdf}. \\\hline

\textbf{Actividad 6 (Sesión 13):} \textit{Trabajo Final · Una habilidad, un PDA, cinco minutos} &
\textbf{Actualizado hoy 20-nov 19:00} con 5 archivos (Video limpio \texttt{4:46}, Guión con YouTube/Drive, Reporte Act 6, Presentación Act 6, Reporte Act 5). Nota previa de las 14:45: $25.00/36.00$. &
Informar a la docente que se actualizó hoy a las 19:00 con el video limpio (YouTube y Drive MP4/SRT) y solicitar revisión; llevar copia en USB por si desea verlo en el aula. &
\texttt{07\_Actividad-6\_.../} (MP4, SRT, \texttt{Guion Trabajo Final\_...docx}, reportes, presentación y \texttt{act6\_video\_32N.ggb}). \\\hline
\end{longtable}}

\section{¿En qué archivo exacto se resuelve cada producto solicitado?}
Para localizar sin pérdida de tiempo cada documento requerido antes de imprimir o pasar a USB:

\begin{enumerate}[leftmargin=1.5em,itemsep=4pt]
  \item \textbf{Examen Parcial 1 (\textit{Parcial 1. Sospechoso, muy sospechoso}):}\\
  Archivo: \texttt{04\_Examen-Parcial-1\_Sospechoso-muy-sospechoso/}\\
  \hspace*{1.2em}\texttt{Parcial-1\_Sospechoso-muy-sospechoso\_Martin-Jonathan-de-la-Cruz-Munoz.pdf} (y editable \texttt{.docx}).\\
  \textit{Clave de respuestas verificadas:} \textbf{1. Cierta} (densidad decimal: $m=\frac{a+b}{2}$); \textbf{2. Falsa} (invariancia aditiva: $(p+r)-(q+r)=p-q$); \textbf{3. Cierta} (propiedad uniforme aditiva reversible); \textbf{4. Falsa} (la mediatriz es recta infinita, solo el punto medio pertenece a $\overline{AB}$); \textbf{5. Cierta} (criterio AAL en $\triangle OHP \cong \triangle OKP \implies PH=PK$).

  \item \textbf{Guion de la Habilidad (Actividad 5, hoja verde con encabezado horizontal):}\\
  Archivo: \textbf{Página 1} de \texttt{06\_Actividad-5\_Habilidades-y-Retos-Sesion-12/}\\
  \hspace*{1.2em}\texttt{Actividad-05\_Guion-Habilidad-y-Retos-Resueltos\_Martin-Jonathan-de-la-Cruz-Munoz.pdf} (complementado con los 6 apartados del Trabajo Final en \texttt{Guion 5\_...docx} y \texttt{reporte-...act-05...pdf}). Habilidad eje: \textbf{Generalizar} ($F_2=k\cdot F_1$ con $k=(d_2/d_1)^2$, CONALEP).

  \item \textbf{Retos de la Sesión 12 / Actividad 5 (\textit{Solución al reto}, tabla verde de habilidades):}\\
  Archivo: \textbf{Página 2} de \texttt{06\_Actividad-5\_Habilidades-y-Retos-Sesion-12/}\\
  \hspace*{1.2em}\texttt{Actividad-05\_Guion-Habilidad-y-Retos-Resueltos\_Martin-Jonathan-de-la-Cruz-Munoz.pdf} y \textbf{Sección 1 + Figura 1} de \texttt{reporte-...act-05...pdf}. Incluye las 8 habilidades SEP con cálculo y esquema TikZ: \textit{Calcular} ($90$), \textit{Inferir} (total $34$, $\bar{x}=6.8$, pico Jueves $=10$), \textit{Comunicar} ($\frac{3}{4}=0.75$), \textit{Medir} ($2.88\text{ m}^2$), \textit{Imaginar} (rotación horaria $90^\circ$ «L» $\to$ «$\Gamma$»), \textit{Estimar} ($\approx 11\text{--}12\text{ cm}$), \textit{Generalizar} ($a_n=n^2$) y \textit{Deducir} ($\angle B=\angle C=70^\circ$).

  \item \textbf{Retos de Actividad 2 (\textit{El lugar perfecto}):}\\
  \hspace*{1.2em}\texttt{02\_Actividad-2\_El-lugar-perfecto/}\\
  \hspace*{1.2em}\texttt{Retos-Resueltos-Actividad-02\_El-Lugar-Perfecto\_Martin-Jonathan-de-la-Cruz-Munoz.pdf} (3 pp.) y Sección 1 de \texttt{reporte-...act-02...pdf}.

  \item \textbf{Retos de Actividad 3 (\textit{Universo 3-60}):}\\
  \hspace*{1.2em}\texttt{03\_Actividad-3\_Universo-3-60/}\\
  \hspace*{1.2em}\texttt{Retos-Resueltos-Actividad-03\_Universo-3-60\_Martin-Jonathan-de-la-Cruz-Munoz.pdf} (4 pp.) y Secciones 1--2 de \texttt{reporte-...act-03...pdf}.

  \item \textbf{Portafolio Bloque I:} \texttt{00\_Portafolio-Evidencias-Bloque-I/}\\
  \hspace*{1.2em}\texttt{Portafolio-Evidencias-Bloque-I\_Completo\_Martin-Jonathan-de-la-Cruz-Munoz.pdf} (39 pp.: portada con índice paginado p. 1 + Act 1 fracciones/decimales pp. 2--10, Act 2 suma/resta pp. 11--17, Act 3 multiplicación/signos pp. 18--27, Act 4 proporcionalidad pp. 28--33, Act 5 ecuaciones lineales pp. 34--39).
\end{enumerate}

\section{Guión práctico y checklist para mañana con la Mtra. Denis}

\begin{center}
\begin{tikzpicture}[scale=0.84,font=\scriptsize,node distance=0.45cm]
  \node[draw=iiiepeNavy,fill=iiiepeNavy!10,rounded corners=3pt,text width=4.1cm,align=center,inner sep=5pt] (p1)
    {\textbf{1. Llevar impreso / físico}\\Portafolio Bloque I (39 pp.)\\Parcial 1 resuelto\\Hojas de Retos Act 2, Act 3 y Act 5};
  \node[draw=iiiepeBlue,fill=iiiepePale,rounded corners=3pt,text width=4.4cm,align=center,inner sep=5pt,right=0.6cm of p1] (p2)
    {\textbf{2. Respaldo en USB}\\Carpetas \texttt{00\_} a \texttt{07\_}\\(Omitir \texttt{\_respaldo\_sistema/})\\Archivos \texttt{.ggb} + Video MP4/SRT};
  \node[draw=iiiepeAqua!80!black,fill=iiiepeAqua!12,rounded corners=3pt,text width=4.8cm,align=center,inner sep=5pt,right=0.6cm of p2] (p3)
    {\textbf{3. Aclarar en EVA / Plataforma}\\Trabajo Final actualizado hoy 19:00\\(YouTube: \texttt{youtu.be/xZOLQRSTyQE})\\Consultar si sube Portafolio Bloque I};
  \draw[-{Latex[length=2mm]},thick,iiiepeNavy] (p1) -- (p2);
  \draw[-{Latex[length=2mm]},thick,iiiepeNavy] (p2) -- (p3);
\end{tikzpicture}
\end{center}

\textbf{Puntos concretos para conversar con la profesora mañana:}
\begin{enumerate}[leftmargin=1.5em,itemsep=2pt]
  \item \textbf{Portafolio de Evidencias Bloque I:} Entregar el tomo completo de 39 páginas (o en USB) y preguntar si además de recibirlo presencialmente requiere que se suba el mismo PDF (\texttt{1.43 MB}) al buzón abierto en plataforma.
  \item \textbf{Actividades presenciales sin buzón:} Entregar en físico el \textit{Parcial 1} resuelto, las hojas de retos de \textit{Actividad 2} y \textit{Actividad 3}, las 2 hojas verdes de \textit{Actividad 5} (Guion de la Habilidad y Solución a los retos) y confirmar la calificación de \textit{Actividad 1} en el tablero Canva del \textbf{Equipo 4}.
  \item \textbf{Trabajo Final actualizado hoy a las 19:00:} Comentar que en plataforma quedó cargada hoy la versión corregida a las \textbf{19:00} (video limpio de \textbf{4:46 min} en \url{https://youtu.be/xZOLQRSTyQE} y respaldo en Drive con subtítulos \texttt{.srt} independientes, corrigiendo a $8\text{ N} \to 32\text{ N}$ y alineando las técnicas didácticas UnADM), para solicitar la actualización de la calificación registrada a las 14:45 ($25/36$).
  \item \textbf{Nota de prevención técnica:} Si copias las entregas a una memoria USB para mostrarlas o entregarlas en digital, copia únicamente las carpetas numeradas \texttt{00\_} a \texttt{07\_} y este informe guía; la carpeta \texttt{\_respaldo\_sistema/} es archivo interno de respaldo histórico y no debe mezclarse en la entrega.
\end{enumerate}

\end{document}
